\chapter{Installing, Verifying, Upgrading, Removing}
\label{ch:install}

\section{Install}
\ifarch
\begin{lstlisting}[style=console]
$ sudo pacman -U kernel-borderlands-0.1.0-1-x86_64.pkg.tar.zst
resolving dependencies...
looking for conflicting packages...
Packages (1) kernel-borderlands-0.1.0-1
:: Proceed with installation? [Y/n]
(1/1) installing kernel-borderlands
group kb already exists, skipping      # or: created group: kb
created user: operator (member of kb)
generated SSH host key: /etc/kb/ssh/ssh_host_ed25519_key
created empty /etc/kb/authorized_keys - add operator public keys here
provisioning complete
==> Kernel Borderlands installed. Nothing has been started.
\end{lstlisting}
\else
\begin{lstlisting}[style=console]
$ sudo apt install ./kernel-borderlands_0.1.0_amd64.deb
Reading package lists... Done
The following NEW packages will be installed:
  kernel-borderlands
Setting up kernel-borderlands (0.1.0) ...
created group: kb
created user: operator (member of kb)
generated SSH host key: /etc/kb/ssh/ssh_host_ed25519_key
created empty /etc/kb/authorized_keys - add operator public keys here
provisioning complete
\end{lstlisting}
Use \texttt{./} in front of the filename: without it \cmd{apt} looks for a package \emph{name} in the archive.
\fi

\subsection{What the post-install step did}
\begin{enumerate}
\item Created group \texttt{kb} and system user \texttt{operator} (no home, shell \file{/usr/sbin/nologin}) if absent.
\item Created \file{/etc/kb} (\texttt{root:kb}, 0750) and \file{/etc/kb/ssh/}.
\item Generated \file{/etc/kb/ssh/ssh_host_ed25519_key} \emph{only if it did not exist}. Re-installing never rotates it.
\item Created an empty \file{/etc/kb/authorized_keys} (\texttt{operator:kb}, 0600).
\item Ran \texttt{systemctl daemon-reload} so the new units are visible.
\item Ran \file{/usr/lib/kb/migrate-legacy.sh} (Section~\ref{sec:migrate}): repaired a legacy operator-sshd install if one exists, and removed any boot autostart of \cmd{kb-checker}.
\end{enumerate}
It did \textbf{not} enable or start anything (on a fresh install it also made sure \cmd{kb-checker} is stopped).

\section{Automatic repair of a legacy dev install}
\label{sec:migrate}
If the host previously ran the development install (\texttt{install.sh} or \texttt{make install}), the package
repairs it for you; there is nothing to do by hand. \file{/usr/lib/kb/migrate-legacy.sh} runs from the post-install
step, is idempotent, and does the following only when it finds the legacy layout:
\begin{enumerate}
\item Moves \file{/etc/ssh/sshd_config.d/kb-operator.conf} to \file{/etc/kb/sshd_kb-operator.conf.legacy}. It is kept as a backup, never deleted, and your edits are not lost; the \emph{packaged} \file{sshd_kb-operator.conf} stays in force because its \texttt{ForceCommand} points at the packaged wrapper.
\item Repoints a legacy \file{/etc/systemd/system/sshd@kb-operator.service} at \file{/etc/kb/sshd_kb-operator.conf} (backup: \file{.pre-migration}).
\item Runs \texttt{systemctl daemon-reload} and \emph{reloads} the host's own \cmd{sshd} (\unit{ssh} on Ubuntu, \unit{sshd} on Arch), which releases the port 2222 it had wrongly bound. Reload, not restart: your admin session stays up.
\item Restarts \unit{sshd@kb-operator} \textbf{unless you stopped it by hand} (it is neither enabled nor active); in that case it is left exactly as you left it.
\end{enumerate}
Separately, on every install and upgrade it removes any boot autostart of \cmd{kb-checker} (and stops it on a fresh
install), so the watchdog stays dead unless you start it yourself (Chapter~\ref{ch:checker}). Units from a legacy
install that live in \file{/etc/systemd/system} still override the packaged ones; the shadowing note at the end of this chapter applies.

\begin{kbwarn}[THE OPERATOR ACCOUNT USES A NOLOGIN SHELL: VERIFY THE CONSOLE ON YOUR HOST]
\file{provision.sh} creates \texttt{operator} with \file{/usr/sbin/nologin}. OpenSSH runs a \texttt{ForceCommand}
as \texttt{\$SHELL -c ...}; on many systems a \texttt{nologin} shell therefore refuses it. This has
\textbf{not} been tested with the packaged account because the build sessions had no root. Chapter~\ref{ch:ssh}
gives the one-command test and the fix (give \texttt{operator} the shell \file{/bin/sh}; the sshd config
already forbids forwarding and forces the wrapper).
\end{kbwarn}

\section{Verify the installation}
\resetsteps
\step{The files are where they should be}
\ifarch
\begin{lstlisting}[style=console]
$ pacman -Ql kernel-borderlands | grep -E '/usr/bin/|/etc/kb/'
\end{lstlisting}
\else
\begin{lstlisting}[style=console]
$ dpkg -L kernel-borderlands | grep -E '/usr/bin/|/etc/kb/'
$ dpkg -s kernel-borderlands | grep -E 'Status|Version'
Status: install ok installed
\end{lstlisting}
\fi

\step{The units are known to systemd and use the right paths}
\begin{lstlisting}[style=console]
$ systemctl cat kbd.service | grep -E 'ExecStart|Environment'
ExecStart=/usr/bin/kbd --db /var/lib/kbd/state.db --policy /etc/kb/policy.yaml \
          --rules /etc/kb/rules.yaml --workloads /etc/kb/workloads.yaml
Environment=KB_HTTP_CORS_ORIGIN=http://127.0.0.1:8082,http://localhost:8082
$ systemd-analyze verify kbd.service kb-sensor.service kbopd.service
\end{lstlisting}

\step{The CLI answers and reports the maintenance hold}
\begin{lstlisting}[style=console]
$ kbctl status
kbd.service                  inactive
kb-sensor.service            inactive
kbopd.service                inactive
sshd@kb-operator.service     inactive
kb-checker.service           UNDER MAINTENANCE AND REPAIRS (kept down)
\end{lstlisting}

\step{The provisioning result}
\begin{lstlisting}[style=console]
$ getent passwd operator
operator:x:998:994::/home/operator:/usr/sbin/nologin
$ sudo ls -l /etc/kb /etc/kb/ssh
\end{lstlisting}

\section{Upgrade}
\ifarch
Install the new file the same way (\texttt{pacman -U}). Because the five config files are in the
\texttt{backup=()} array, a locally modified file is \emph{kept} and the new default is written beside it as
\file{.pacnew}. Merge with \texttt{pacdiff} (from \texttt{pacman-contrib}).
\else
Install the new file the same way (\texttt{apt install ./new.deb}). The five files in \file{conffiles} are
handled by \cmd{dpkg}: if you changed one and the package changed it too, you are asked whether to keep yours or
take the maintainer's; answer \texttt{N} (keep yours), then diff against the \file{.dpkg-dist} copy.
\fi
Restart what changed: \texttt{sudo systemctl daemon-reload \&\& sudo kbctl down \&\& sudo kbctl up}.
The SQLite state in \file{/var/lib/kbd/} is untouched by an upgrade.

\section{Remove}
\ifarch
\begin{lstlisting}[style=console]
$ sudo pacman -R kernel-borderlands
\end{lstlisting}
\else
\begin{lstlisting}[style=console]
$ sudo apt remove kernel-borderlands      # keeps modified config
$ sudo apt purge  kernel-borderlands      # removes config files the package owns
\end{lstlisting}
\fi
Before removal the units are disabled and stopped. These are \textbf{intentionally left behind}, even after a
purge, because they hold credentials and audit state an operator would not expect a package manager to destroy:
\file{/etc/kb/ssh/} and \file{/etc/kb/authorized_keys}, \file{/var/lib/kbd/state.db}, the \texttt{kb} group and the
\texttt{operator} user. To remove them as well:
\begin{lstlisting}[style=console]
$ sudo rm -rf /etc/kb /var/lib/kbd
$ sudo userdel operator && sudo groupdel kb
\end{lstlisting}

\begin{kbnote}[IF A \emph{DEVELOPMENT} INSTALL ALREADY EXISTS ON THE HOST]
\texttt{install.sh} and \texttt{sudo make install} place binaries in \file{/usr/local/bin} and units in
\file{/etc/systemd/system}. Units in \file{/etc/systemd/system} \textbf{override} the package's units in
\file{/usr/lib/systemd/system}, and \file{/usr/local/bin} wins over \file{/usr/bin} on \texttt{PATH}. Installing the
package on top of such a host will silently keep running the old code. Remove the old units and binaries
first (stop them, delete the files, \texttt{daemon-reload}), then install the package.
\end{kbnote}
